\documentclass[12pt]{article}
\setlength{\topmargin}{-0.5 in}
\setlength{\textheight}{9.5 in}
\setlength{\oddsidemargin}{-0.4 in}
\setlength{\textwidth}{7.0 in}
\usepackage{amsmath,graphicx,amssymb, diagbox, braket}
\usepackage{tikz}
\usepackage{braket}
\newcommand*\circled[1]{\tikz[baseline=(char.base)]{
		\node[shape=circle,draw,inner sep=2pt] (char) {#1};}}
\pagestyle{empty}
\begin{document}
	
\centerline{QC HW 8
	\hspace{0.4 in}V. Kamat
	\hspace{0.4 in}Spring 2026}
	\begin{quote}
		{\sf
		{\bf Directions.} For Questions 1 and 2, refer to the Week 7 notes (on Grover search) for additional context surrounding each of the exercises below.  
		}
	\end{quote}
\begin{enumerate}
\item 
\begin{enumerate}
\item Show that the outer product $\ket{s}\bra{s}=\frac{1}{N}J_N$, where $N=2^n$, $J$ is the $N\times N$ all $1$'s matrix, and $\ket{s}$ is the uniform superposition over all $n$-qubit computational basis states.
\item \label{refd} Show that for the diffusion operator given by $D=2\ket{s}\bra{s}-I_N$, $D\ket{s}=\ket{s}$, and $D\ket{s^{\perp}}=-\ket{s^{\perp}}$ for any $\ket{s^{\perp}}$ that is orthogonal to $\ket{s}$.
\item As defined in the notes, let $\ket{g}$ be the unique marked basis state, and let $\ket{b}$ be the uniform superposition over all of the remaining $N-1$ unmarked basis states. Show that $\braket{g|b}=0$, i.e. $\ket{g}$ and $\ket{b}$ are orthogonal.
\item \label{refz} Show that $Z_f\ket{g}=-\ket{g}$ and $Z_f\ket{b}=\ket{b}$, where $Z_f$ is the unitary phase oracle implementing the mapping $\ket{x}\to (-1)^{f(x)}\ket{x}$. 
\end{enumerate}
\item Rework the analysis of Grover's algorithm from the notes with the assumption that there are $1\leq M<< N$ marked states, and show that the optimal number of Grover iterations is $k\approx \dfrac{\pi}{4}\sqrt{\dfrac{N}{M}}-\dfrac{1}{2}.$
\item Show that if $U$ is a unitary matrix with eigenvector $\ket{\psi}$ such that $U\ket{\psi}=\lambda\ket{\psi}$, $\lambda$ is a complex number of unit modulus (magnitude). 
\end{enumerate}
\end{document}